\par\medskip\noindent
\underline{Definición y Proposición 2.11.} Un cogrupoide
$(R_0,R_1,\ldots)$ en $\widehat{\underline C}_\Lambda$ se llama
\underline{liso} si se cumple una, y por tanto ambas, de las siguientes
condiciones equivalentes:

(i) $R_1$ es un anillo de series de potencias formales sobre $R_0$ por
medio de $s_1$ (y por medio de $s_2$);

(ii) el funtor
\[
 \epsilon:\underline h_{R_0}
 \longrightarrow\underline h_{(R_0,R_1,\ldots)}
\]
es liso sobre $\underline C_\Lambda$.

Un \underline{cogrupoide mínimamente liso}, llamado también
\underline{cogrupoide liso normalizado}, en
$\widehat{\underline C}_\Lambda$ es un cogrupoide liso
$(R_0,R_1,\ldots)$ para el cual
$\underline h_{(R_0,R_1,\ldots)}(k[\varepsilon])$ es totalmente
disconexo; equivalentemente,
\[
 s_1(k[\varepsilon])=s_2(k[\varepsilon]).
\]

\underline{Demostración.} La lisura de
$\epsilon:\underline h_{R_0}\to
\underline h_{(R_0,R_1,\ldots)}$ significa que, para toda aplicación
suprayectiva
\[
 (S',f')\xrightarrow{(h,u)}(S,f)
\]
en $\underline h_{(R_0,R_1,\ldots)}$, existen un levantamiento $f''$ de
$f$ a $S'$ y un isomorfismo
\[
 (S',f'')\xrightarrow{(1,u')}(S',f').
\]
Equivalentemente, dados $u\in h_{R_1}(S)$ y un levantamiento $f'$ de
$s_1(u)$ a $S'$, existe un levantamiento $u'$ de $u$ a $S'$ tal que
$s_1(u')=f'$, para toda suprayección $S'\to S$ en
$\underline C_\Lambda$. Esto es precisamente la lisura de
\[
 s_1:\underline h_{R_1}\longrightarrow\underline h_{R_0}
\]
sobre $\underline C_\Lambda$, y esta, a su vez, afirma que $R_1$ es un
anillo de series de potencias formales sobre $R_0$ por medio de $s_1$.
\hfill$\square$

\underline{Proposición 2.12.} Sea $(R_0,R_1,\ldots)$ un cogrupoide en
$\widehat{\underline C}_\Lambda$.

(a) Si $(R_0,R_1,\ldots)$ es liso, entonces
$\underline h_{(R_0,R_1,\ldots)}$ es homogéneo.

(b) Si $s_1(k[\varepsilon])=s_2(k[\varepsilon])$, entonces las condiciones
siguientes son equivalentes:

(i) $(R_0,R_1,\ldots)$ es liso;

(ii) $\underline h_{(R_0,R_1,\ldots)}$ es homogéneo;

(iii) $\underline h_{(R_0,R_1,\ldots)}$ es semihomogéneo.

%%%%%%%%%%  scan idx 75  |  folio 64  %%%%%%%%%%

(c) Si $\underline h_{(R_0,R_1,\ldots)}$ es semihomogéneo, entonces
es equivalente a $\underline h_{(S_0,S_1,\ldots)}$ para algún cogrupoide
mínimamente liso $S_*$.

\underline{Demostración.} (a) Considérese un diagrama
\[
\begin{tikzcd}[column sep=huge,row sep=large]
 (S_1,f_1) \arrow[dr,"{(h_1,u_1)}"']
 && (S_2,f_2) \arrow[dl,"{(h_2,u_2)}"]\\
 & (S,f) &
\end{tikzcd}
\]
en $\underline h_{(R_0,R_1,\ldots)}$, donde
$h_2:S_2\longrightarrow S$ es suprayectiva en
$\underline C_\Lambda$. Puesto que
$s_1:\underline h_{R_1}\longrightarrow\underline h_{R_0}$ es
liso, $u_1^{-1}u_2$ se levanta a un elemento $v\in h_{R_1}(S_2)$ con
$s_1(v)=f_2$. Se tiene
\[
 h_2\circ s_2(v)=h_1\circ f_1,
\]
por lo que $f_1\times s_2(v)$ es un elemento de
$h_{R_0}(S_1\times_SS_2)$ y proporciona un diagrama conmutativo
\[
\begin{tikzcd}[column sep=huge,row sep=large]
 & (S_1\times_SS_2,f_1\times s_2(v))
   \arrow[dl,"{\pi_1}"']
   \arrow[dr,"{(\pi_2,v^{-1})}"] &\\
 (S_1,f_1) \arrow[dr,"{(h_1,u_1)}"']
 && (S_2,f_2) \arrow[dl,"{(h_2,u_2)}"]\\
 & (S,f) &
\end{tikzcd}
\]
Afirmamos que su objeto superior es el producto fibrado
\[
 (S_1,f_1)\times_{(S,f)}(S_2,f_2)
\]
en $\underline h_{(R_0,R_1,\ldots)}$. En efecto, para todo objeto
$(T,g)$, una aplicación al objeto superior equivale a una tupla
$(h'_1\times h'_2,w_1,w_2)$ que satisface
\begin{align*}
 h_1w_1&=h_2w_2,&
 s_1(w_1)&=h'_1g,&s_2(w_1)&=f_1,\\
 s_1(w_2)&=h'_2g,&
 s_2(w_2)&=s_2(v).
\end{align*}
Al poner $w'_2=v^{-1}w_2$ se obtiene una biyección con las tuplas
$(h'_1\times h'_2,w_1,w'_2)$ que satisfacen
\begin{align*}
 u_2h_2w'_2&=u_1h_1w_1,&
 s_1(w_1)&=h'_1g,&s_2(w_1)&=f_1,\\
 s_1(w'_2)&=h'_2g,&
 s_2(w'_2)&=f_2.
\end{align*}
Estas son exactamente las parejas de flechas de $(T,g)$ a $(S_1,f_1)$
y $(S_2,f_2)$ cuyos compuestos hacia $(S,f)$ coinciden. Por tanto,
\begin{multline*}
 \operatorname{Hom}((T,g),(S_1\times_SS_2,
                 f_1\times s_2(v)))\\
 \simeq
 \operatorname{Hom}((T,g),(S_1,f_1))
 \times_{\operatorname{Hom}((T,g),(S,f))}
 \operatorname{Hom}((T,g),(S_2,f_2)),
\end{multline*}
lo que demuestra (a).

%%%%%%%%%%  scan idx 76  |  folio 65  %%%%%%%%%%

(b) Solo queda demostrar (iii)$\Rightarrow$(i). Si
$\underline h_{(R_0,R_1,\ldots)}$ es semihomogéneo, entonces
\[
 \dim[\underline h_{(R_0,R_1,\ldots)}(k[\varepsilon])]
 =\dim h_{R_0}(k[\varepsilon])<\infty,
\]
por lo que admite un objeto formal mínimamente versal. Por otra parte,
$(R_0,1_{R_0})$ es un objeto mínimo cuasi-inicial de
$\widehat{\underline h}_{(R_0,R_1,\ldots)}$, porque
\[
 h_{R_0}(k[\varepsilon])
 \longrightarrow
 [\underline h_{(R_0,R_1,\ldots)}(k[\varepsilon])]
\]
es biyectiva. Por tanto, la Proposición 1.10(ii) hace de
$(R_0,1_{R_0})$ un objeto formal mínimamente versal de
$\underline h_{(R_0,R_1,\ldots)}$. Equivalentemente,
\[
 \epsilon:\underline h_{R_0}
 \longrightarrow\underline h_{(R_0,R_1,\ldots)}
\]
es mínimamente liso. Esto demuestra (i).
